\documentclass[12pt]{article}
\usepackage{amsmath,amsfonts,fullpage,amsthm}

\newtheorem*{theorem}{Theorem}
\newtheorem*{fact}{Fact}
\newtheorem*{goal}{Goal}
\newtheorem*{bigidea}{Big Idea}
\newtheorem*{idea}{Idea}
\newtheorem*{defn}{Definition}
\newtheorem*{ex}{Example}

%Style section
%\setlength{\textheight}{10in}
%\setlength{\textwidth}{7in}
%\hoffset=-0.75in
 \pagestyle{plain}
% \setlength{\topmargin}{0in}
% \setlength{\headsep}{0in}
% \setlength{\oddsidemargin}{0in}
% \setlength{\evensidemargin}{0in}

\begin{document}


\thispagestyle{empty} %use this to get rid of page numbers

\noindent
{Math 166 \hskip 2.5 truein  Names:\ \hrulefill}

%\vskip 0.3truein
%\bigskip

\noindent
{Fun with Sequences }

\begin{bigidea} \rm
By the end of Chapter 11, we will be able to approximate \emph{any function} by a polynomial, called its \emph{Taylor polynomial}. Polynomials are easy to integrate, so this will help us to evaluate difficult/impossible integrals.
\end{bigidea}
%\medskip



\begin{defn} \rm
We say a sequence $\{a_n\}$ \textbf{converges} to $L$ if $a_n\rightarrow L$ as $n\rightarrow\infty$, that is, if $\displaystyle\lim_{n\rightarrow\infty} a_n=L$. If the limit $\displaystyle\lim_{n\rightarrow\infty} a_n$ does not exist, we say $\{a_n\}$ \textbf{diverges}.
%(Note: This definition is a bit informal. For the formal definition of convergence and divergence of a sequence, see p.\ $514$.)
\end{defn}

\begin{theorem}
If $\displaystyle\lim_{x\rightarrow\infty}f(x)$ exists, then the sequence $a_n=f(n)$ converges to the same limit: \[\displaystyle\lim_{n\rightarrow\infty}a_n=\displaystyle\lim_{x\rightarrow\infty}f(x)\]
\end{theorem}

\begin{enumerate}
  \setcounter{enumi}{0}
\item Let $a_n= r^n$ for some real number $r$ (so $\{a_n\}$ is a \textbf{geometric sequence}). What is the limit of the sequence $\{a_n\}$ for the following values of $r$? Can you say in general for which values of $r$ the sequence converges?
\begin{enumerate}
\item $r=5$ \vspace{.5in}
\item $r=1$\vspace{.5in}
\item $r=\frac{1}{2}$\vspace{.5in}
\item $r=0$\vspace{.5in}
\item $r=-\frac{1}{3}$\vspace{.5in}
%\item $r=-1$\vspace{.2in}
\item $r=-4$\vspace{.5in}
\end{enumerate}

\medskip
\noindent
The following theorem says we may ``bring a limit of a sequence inside a continuous function.''
\begin{theorem}
If $f$ is continuous and $\displaystyle\lim_{n\rightarrow\infty} a_n=L$, then
\[\displaystyle\lim_{n\rightarrow\infty} f(a_n) = f\left(\displaystyle\lim_{n\rightarrow\infty} a_n\right)=f(L)\]
\end{theorem}

\item Use this theorem to find the limit of the following sequences:
Start with $n=1$
\begin{enumerate}
\item $a_n=\cos\left(\displaystyle\frac{\pi n^4}{3+n^4}\right)$

%\smallskip
%\noindent
%$\displaystyle\lim_{n\rightarrow\infty}\cos\left(\displaystyle\frac{\pi n^4}{3+n^4}\right)=\cos\left(\displaystyle\lim_{n\rightarrow\infty}\displaystyle\frac{\pi n^4}{3+n^4}\right)=$
\vspace{.5in}
\bigskip
\item $b_n=\sqrt{5+\frac{2}{n}}$
\vspace{.5in}
\bigskip
\item $c_n=e^{\frac{4n}{n^2+1}}$
\vspace{.5in}
\end{enumerate}
\end{enumerate} 
\begin{defn} \rm
%\begin{itemize}
%\item
A sequence $\{a_n\}$ is \emph{bounded above} by $M$ if $a_n\leq M$ for all $n$.
%\item
$\{a_n\}$ is \emph{bounded below} by $m$ if $a_n\geq m$ for all $n$.
%\item
$\{a_n\}$ is \emph{bounded} if it is both bounded below and bounded above.
%\end{itemize}
\end{defn}

\begin{enumerate}
\setcounter{enumi}{2}
\item Classify the sequences in Problem 1 as bounded above, bounded below, bounded, or none of these.
\end{enumerate}

\bigskip

\pagebreak

\begin{defn} \rm
%\begin{itemize}
%\item
A sequence $\{a_n\}$ is \emph{increasing} if $a_n\leq a_{n+1}$ for all $n$.
%\item
$\{a_n\}$ is \emph{decreasing} if $a_n\geq a_{n+1}$ for all $n$.
%\item
 $\{a_n\}$ is \emph{monotonic} if it is either increasing or decreasing.


%\end{itemize}
\end{defn}

\begin{enumerate}
\setcounter{enumi}{3}
\item Classify the sequences in Problem 1 as increasing, decreasing, or neither.


\vfill

\item If a sequence $\{a_n\}$ converges, does it have to be bounded?

\vfill

\item If a sequence $\{a_n\}$ is bounded, does it have to converge?

\vfill

\item If a sequence $\{a_n\}$ is bounded and monotonic, does it have to converge?
\vfill
\thispagestyle{empty}
\end{enumerate}



\end{document}
